%

\pol{\section{Origami}}
\ita{\section{Origami}}
\label{section_origami}
\index{origami}

\pol{Aksjomaty Huzity-Justina (albo Huzity-Hatoriego) to zestaw siedmiu aksjomatów, które opisują wszystkie możliwe pojedyncze zgięcia w geometrii origami na płaszczyźnie.}
\ita{Gli assiomi di Huzita-Justin (o di Huzita-Hatori) sono un insieme di sette assiomi che descrivono tutte le possibili pieghe singole nella geometria origami sul piano.}
\pol{\index{aksjomat!Huzity-Hatoriego}}\ita{\index{assioma!di Huzita-Hatori}}%
\pol{\index{aksjomat!Huzity-Justina}}\ita{\index{assioma!di Huzita-Justin}}%
\pol{Niechf $p_1, p_2$ będą dwoma (różnymi) punktami, zaś $l_1, l_2$ dwoma różnymi prostymi.}
\ita{Siano $p_1, p_2$ due punti (distinti), e $l_1, l_2$ due rette distinte.}

\begin{enumerate}
    \item \pol{Istnieje dokładnie jedno zgięcie, które przechodzi przez punkty $p_1$ i $p_2$.}
    \ita{Esiste esattamente una piega che passa per i punti $p_1$ e $p_2$.}
    \item \pol{Istnieje dokładnie jedno zgięcie, które przenosi punkt $p_1$ na $p_2$.}
    \ita{Esiste esattamente una piega che porta il punto $p_1$ su $p_2$.}
    \item \pol{Istnieje zgięcie, które przenosi prostą $l_1$ na $l_2$.}
    \ita{Esiste una piega che porta la retta $l_1$ su $l_2$.}
    \item \pol{Istnieje dokładnie jedno zgięcie przez $p_1$, które jest prostopadłe do $l_1$.}
    \ita{Esiste esattamente una piega passante per $p_1$ che è perpendicolare a $l_1$.}
    \item \pol{Istnieje zgięcie przez $p_2$, które przenosi punkt $p_1$ na prostą $p_2$.}
    \ita{Esiste una piega passante per $p_2$ che porta il punto $p_1$ sulla retta $p_2$.}
    \item \pol{Istnieje zgięcie, które przenosi $p_1$ na $l_1$, zaś $p_2$ na $l_2$.}
    \ita{Esiste una piega che porta $p_1$ su $l_1$ e $p_2$ su $l_2$.}
    \item \pol{Istnieje zgięcie prostopadłe do $l_2$, które przenosi $p_1$ na $l_2$.}
    \ita{Esiste una piega perpendicolare a $l_2$ che porta $p_1$ su $l_2$.}
\end{enumerate}

\pol{Zestaw nie jest minimalny, ale wyczerpuje wszystkie możliwe ruchy.}
\ita{L'insieme non è minimale, ma esaurisce tutti i movimenti possibili.}
\pol{Wszystkie znajdzie francuski składacz Jacques Justin \cite{justin_1986} w 1986 roku.}
\ita{Tutti saranno trovati dal piegatore francese Jacques Justin \cite{justin_1986} nel 1986.}
\index[persons]{Justin, Jacques}%
\pol{Później część z nich będzie odkrywana na nowo: przez japońsko-włoskiego matematyka Humiakiego Huzitę (od 1 do 6, około 1991 roku), przez Davida Aucklyego i Johna Clevelanda \cite{cleveland_1995} (od 1 do 5 w 1995 roku), Koshiro Hatoriego (aksjomat nr 7 w 2001 roku) i Roberta Langa \cite{alperin_2009} (także aksjomat nr 7).}
\ita{In seguito alcuni di essi saranno riscoperti: dal matematico nippo-italiano Humiaki Huzita (dall'1 al 6, intorno al 1991), da David Auckly e John Cleveland \cite{cleveland_1995} (dall'1 al 5 nel 1995), da Koshiro Hatori (assioma n. 7 nel 2001) e da Robert Lang \cite{alperin_2009} (anch'egli l'assioma n. 7).}
\index[persons]{Huzita, Humiaki}%
\index[persons]{Auckly, David}%
\index[persons]{Cleveland, John}%
\index[persons]{Hatori, Koshiro}%
\index[persons]{Lang, Robert}%
% TODO: Alperin =  Roger C. Alperin; Robert J. Lang (2009). "One-, Two-, and Multi-Fold Origami Axioms" (PDF). 4OSME. A K Peters. Archived from the original (PDF) on 2022-02-13. Retrieved 2012-04-20.

\pol{Aksjomat nr 2 znajduje symetralną odcinka $p_1 p_2$; nr 3 znajduje dwusieczną kąta, na ramionach którego leżą proste $l_1$, $l_2$.}
\ita{L'assioma n. 2 trova l'asse del segmento $p_1 p_2$; il n. 3 trova la bisettrice dell'angolo sui cui lati giacciono le rette $l_1$, $l_2$.}
\pol{Może to być nieoczywiste, ale aksjomat nr 5 jest równoważny przecięciu okręgu prostej.}
\ita{Può non essere evidente, ma l'assioma n. 5 è equivalente all'intersezione di un cerchio e una retta.}
\pol{Aksjomat nr 6 znajduje prostą styczną do dwóch parabol.}
\ita{L'assioma n. 6 trova una retta tangente a due parabole.}
\pol{(To zgięcie nazywa się zgięciem Beloch, ponieważ Margharita Beloch pokaże w 1936, że jest to równoważne co rozwiązaniu wielomianowych równań trzeciego stopnia).}
\ita{(Questa piega si chiama piega di Beloch, perché Margharita Beloch mostrerà nel 1936 che essa è equivalente alla risoluzione di equazioni polinomiali di terzo grado).}
\index[persons]{Beloch, Margharita}%
% https://en.wikipedia.org/wiki/Mathematics_of_paper_folding

\pol{Aksjomaty od pierwszego do czwartego są słabsze niż cyrkiel i linijka: pozwalają zbudować liczby pitagorejskie, najmniejsze ciało zawierające wszystkie liczby wymierne, które oprócz (dowolnej) liczby $x$ zawierają też $\sqrt{x^2 + 1}$.}
\ita{Gli assiomi dal primo al quarto sono più deboli di riga e compasso: permettono di costruire i numeri pitagorici, il più piccolo campo contenente tutti i numeri razionali che, oltre a un numero (qualunque) $x$, contiene anche $\sqrt{x^2 + 1}$.}
\pol{Aksjomaty od 1 do 5 są dokładnie tak samo mocne jak cyrkiel i linijka, nr 6 pozwala na rozwiązanie problemów delfickiej wyroczni, zaś nr 7 nie daje nic nowego.}
\ita{Gli assiomi da 1 a 5 sono esattamente potenti quanto riga e compasso, il n. 6 permette di risolvere i problemi dell'oracolo di Delfi, mentre il n. 7 non dà nulla di nuovo.}
\pol{Ale zgięcie, które opisuje, jest możliwe do wykonania na kartce papieru, więc aksjomat nie jest całkiem niepotrzebny.}
\ita{Ma la piega che descrive è possibile da eseguire su un foglio di carta, quindi l'assioma non è del tutto inutile.}

\pol{Jorge Lucero w 2017 dołoży ósmy aksjomat, że istnieje zgięcie wzdłuż prostej $l_1$.}
\ita{Jorge Lucero nel 2017 aggiungerà un ottavo assioma, secondo cui esiste una piega lungo la retta $l_1$.}
\index[persons]{Lucero, Jorge}%

\pol{Osiem zadań z konkursu matematycznego origami ,,Żuraw'' znajdziemy w $\Delta_{12}^5$.}
\ita{Otto problemi del concorso matematico di origami ,,Żuraw'' si trovano in $\Delta_{12}^5$.}
\pol{Dwie studentki, Barbara Ciesielska i Agnieszka Kowalczyk, opiszą w $\Delta_{16}^7$ konstrukcję odcinka długości $\sqrt[3]{2}$ oraz trysekcję kąta przy użyciu origami.}
\ita{Due studentesse, Barbara Ciesielska e Agnieszka Kowalczyk, descriveranno in $\Delta_{16}^7$ la costruzione di un segmento di lunghezza $\sqrt[3]{2}$ e la trisezione dell'angolo mediante origami.}

\pol{Todo: Hartshorne \cite[s. 249]{hartshorne2000}.}
\ita{Todo: Hartshorne \cite[p. 249]{hartshorne2000}.}

%